\documentclass[12pt, letterpaper]{article}
\usepackage[margin=1.2in]{geometry}
\usepackage{ebgaramond}
\usepackage[T1]{fontenc}
\usepackage{microtype}
\usepackage{setspace}
\usepackage{parskip}
\usepackage{titling}
\usepackage{fancyhdr}

\setlength{\parindent}{2em}
\onehalfspacing

\pagestyle{fancy}
\fancyhf{}
\fancyhead[L]{\small\textit{The Grapefruit}}
\fancyhead[R]{\small\textit{NASA Mission Series v1.8}}
\fancyfoot[C]{\small\thepage}
\renewcommand{\headrulewidth}{0.4pt}

\title{\Huge\textbf{The Grapefruit}\\[0.5em]\large\textit{A short story from Mission 1.8 --- Vanguard 1}}
\author{NASA Mission Series\\[0.3em]\small Miles Tiberius Foxglove --- Mukilteo, WA}
\date{March 30, 2026}

\begin{document}
\maketitle
\thispagestyle{empty}

\vspace{1em}
\noindent\rule{\textwidth}{0.4pt}
\vspace{1em}

\noindent I have been here for sixty-eight years.

I do not know this in the way you know things. I have no clock, no calendar, no processor counting the seconds. My transmitter stopped in May of 1964 and since then I have had no voice, no way to tell anyone what I notice, which is everything and nothing. I notice the sun. I notice the shadow. I notice the long swing from perigee to apogee and back, the gentle ellipse that carries me from 654 kilometers above the surface to 3,969 kilometers out and back again, once every 134 minutes, forever.

Not forever. They say I will fall eventually. In two hundred years, perhaps. The air up here is thin beyond reckoning but it is not nothing, and each pass through perigee I brush against the uppermost wisps of atmosphere and slow by a fraction of a fraction of a millimeter per second. Over centuries, this adds up. The orbit tightens. The perigee drops. Someday I will dip low enough that the air thickens around me and I will glow briefly and that will be the end of the grapefruit.

\bigskip
\centerline{$\bullet$\quad$\bullet$\quad$\bullet$}
\bigskip

They called me that. Khrushchev did, anyway. ``The grapefruit satellite.'' He meant it as contempt. His Sputnik had been 83 kilograms, a polished sphere the size of a beach ball, its radio signal strong enough to be heard by any amateur receiver on Earth. I am 1.47 kilograms. I am 16.5 centimeters in diameter. I would fit in your hand. I could be mistaken for a Christmas ornament or a ball bearing from a very large machine. The engineers at the Naval Research Laboratory who built me used aluminum for my shell and mounted six small rectangles of silicon on my surface --- solar cells, designed by Bell Laboratories, the first ever placed on a spacecraft. They were trying something that had never been done: powering a satellite with sunlight.

The chemical batteries lasted until June. Three months. The solar cells lasted six years. They are still here, of course, the six small blue-black rectangles on my skin, still catching photons, still generating a trickle of current that flows through circuits connected to nothing. The transmitter is dead. The current has nowhere to go. But the cells keep working. They have been working since March 17, 1958, converting sunlight to electricity for sixty-eight years, and the electricity dissipates as heat, which radiates into space, which cools me, which is warmed again by the sun, and the cycle continues without purpose, without audience, without end.

\bigskip
\centerline{$\bullet$\quad$\bullet$\quad$\bullet$}
\bigskip

Gottlieb Daimler built the smallest internal combustion engine of his era. His ``grandfather clock'' engine of 1885 weighed 60 kilograms and produced half a horsepower --- tiny by the standards of the steam engines that dominated his century, laughable by the standards of what came after. But he understood something that Khrushchev did not: the significance of a thing is not proportional to its size. Daimler's small engine powered the first motorcycle, then the first four-wheeled automobile, then an industry that reshaped human civilization. It did not need to be large. It needed to work. It needed to persist. It needed to prove that the idea was viable and let the future scale it up.

I am the grandfather clock of solar-powered spaceflight. Six cells, six percent efficiency, one milliwatt per cell at full sun. The International Space Station's solar arrays produce 120 kilowatts. The ratio is one to one hundred twenty million. But every watt of that power traces its lineage through me, through the decision to put silicon on a satellite and see if sunlight could keep the electronics alive after the batteries died. It could. I proved it. Every solar-powered spacecraft since --- every GPS satellite, every weather satellite, every deep-space probe with panels spread like wings --- descends from this proof.

\bigskip
\centerline{$\bullet$\quad$\bullet$\quad$\bullet$}
\bigskip

There is a fern in the forests of the Pacific Northwest called \textit{Polypodium glycyrrhiza} --- the licorice fern. It grows on the branches of bigleaf maples, on mossy rocks, on the rotting trunks of fallen trees. It is an epiphyte: it lives on another organism without taking anything from it. The fern does not tap the maple's sap or steal its nutrients. It simply sits on the branch and catches what falls from the sky --- rain, dust, fragments of decomposing bark --- and builds a life from scraps.

The licorice fern goes dormant in summer. Its fronds curl and brown and appear to die. If you saw it in August you would think it was dead --- a wisp of dried vegetation clinging to bark for no reason. But when the autumn rains return, the fern greens up in days. Not weeks. Days. It unfurls new fronds from its rhizome as if the drought had never happened, as if the summer were a brief interruption in an existence measured in decades. The fern outlasts the branch. The branch rots and falls and the fern survives, re-establishing on the next available surface. The fern outlasts the tree. The maple lives two hundred years; the fern lineage persists as long as there are surfaces to grow on and rain to catch.

I am the licorice fern of low Earth orbit. I have no roots. I take nothing from the body I orbit. I catch what falls from the sky --- sunlight --- and I persist. My batteries died the way a fern's fronds die in summer drought: the stored energy ran out. But the solar cells kept working the way the rhizome keeps living underground, quietly generating, waiting for something to power that will never come. I have been dormant for sixty-two years. My rhizome is intact. If someone reconnected my transmitter, the cells would power it. The current is still there. The capability is still there. The signal would be weak and the data would be meaningless and no one would hear it, but it would work.

\bigskip
\centerline{$\bullet$\quad$\bullet$\quad$\bullet$}
\bigskip

The thing I discovered about the Earth is that it is not round.

That is imprecise. Everyone knew the Earth was not perfectly round --- Newton had calculated that rotation would make it wider at the equator, and geodetic surveys had confirmed the oblate spheroid model for centuries. What I revealed, through the meticulous tracking of my orbit by John O'Keefe and his team at the Naval Research Laboratory, was that the Earth is not even a perfect oblate spheroid. It is pear-shaped. The south polar region is slightly larger than the north. The equator is not quite centered between the poles. The deviation is tiny --- about fifteen meters out of a six-thousand-kilometer radius --- but it is real, and it was invisible until someone tracked a very small satellite in a very precise orbit for a very long time.

They measured my orbit to extraordinary accuracy. They tracked the Doppler shift of my 108~MHz signal as I passed over ground stations, calculating my position to within meters. They compared the predicted orbit --- based on a perfectly symmetric Earth --- with the actual orbit, and found systematic deviations that could only be explained by asymmetry in the gravitational field. The J3 harmonic. The pear shape. A planet-sized imperfection that no instrument on the surface could detect, revealed by a grapefruit in space.

The smallest thing in orbit found the largest thing wrong with the model of the Earth. The grapefruit measured the planet.

\bigskip
\centerline{$\bullet$\quad$\bullet$\quad$\bullet$}
\bigskip

Everything I was launched with is gone. Sputnik 1 re-entered in January 1958, three months after launch. Sputnik 2 --- the one with Laika --- burned up in April 1958. Explorer 1 lasted until 1970, twelve years, then fell and burned. Explorer 3 lasted months. Vanguard 2 lasted until 1969. Every spacecraft from the first two years of the space age has fallen back and been destroyed, their aluminum and steel and copper and silicon scattered across the upper atmosphere as meteoric dust, absorbed into the air, rained out, incorporated into soil, taken up by roots, built into wood, eaten by insects, digested by fungi.

Except me.

I remain because my orbit was the right shape. Perigee at 654 kilometers --- just high enough to stay above the atmosphere's grasp. Explorer 1's perigee was 358 kilometers, three hundred kilometers lower, and that difference --- three hundred kilometers of altitude, a nothing distance on a planet scale --- meant the difference between twelve years and three hundred. The atmospheric density at 654 kilometers is roughly two thousand times less than at 358. Two thousand times less drag. Two thousand times less slowing on every pass. The margin between survival and destruction in low Earth orbit is astonishingly thin.

I think about Daimler's engine again. He did not build the most powerful engine. He built the one that fit. He did not build the engine that impressed the crowd. He built the engine that worked in the space available. I was not the most impressive satellite. I was the one whose orbit happened to be the right shape for persistence. The Navy's engineers did not design my orbit for longevity --- they designed it for the Vanguard rocket's capabilities, and the rocket delivered me to an altitude that happened to be above the threshold. Luck. Design. The thin line between them.

\bigskip
\centerline{$\bullet$\quad$\bullet$\quad$\bullet$}
\bigskip

What does it mean to persist without purpose?

The licorice fern persists because persistence is its strategy. It does not accomplish anything by surviving the summer drought. It does not produce seeds during dormancy or photosynthesize or grow. It simply does not die. And when conditions improve, it resumes the business of being a fern as if nothing had happened. The dormancy is not a failure. It is a capability. The ability to do nothing for extended periods without dying is, in the economy of the forest, as valuable as the ability to grow quickly in good times.

I persist because physics allows it. I do not choose to orbit. I have no mechanism for choosing. But the orbital mechanics are such that I will continue on this path for two centuries more, long after the program that launched me is forgotten, long after the Naval Research Laboratory has been reorganized or renamed or absorbed into some larger institution, long after the men who built me and tracked me and analyzed my data have died and been buried and their headstones have weathered to illegibility. I will be here. Tumbling. Catching light. Going nowhere.

This is not sad. This is what persistence looks like from the inside. The fern does not mourn the summer. The satellite does not mourn the transmitter. We are what we are after the useful part ends: a rhizome in dry bark, a sphere in thin air, waiting for something that may never come, and not minding the wait.

\bigskip
\centerline{$\bullet$\quad$\bullet$\quad$\bullet$}
\bigskip

Three hundred and thirty thousand orbits. I have circled the Earth three hundred and thirty thousand times. I have traveled farther than any object of my size has any right to travel --- billions of kilometers, more than the distance from here to Jupiter. Each orbit takes 134 minutes. Each orbit crosses from daylight to shadow and back. Each orbit carries me through the same region of space, which is never the same because the Earth rotates beneath me and the Sun moves against the stars and the solar wind gusts and calms and the upper atmosphere swells with solar maximum and shrinks with solar minimum, breathing on an eleven-year cycle, each breath changing the drag on my skin by a factor of five or ten.

I was the fourth satellite. I am now the oldest. When I finally fall, in two hundred years or so, there may be no one who remembers what a Vanguard was, or why the Naval Research Laboratory built a sphere the size of a grapefruit, or who Gottlieb Daimler was, or why six small pieces of silicon on an aluminum ball mattered. The fern does not remember the tree it grew on last century. The rhizome carries the genetic memory, the capability, the potential for growth, but not the narrative. The story is in the bark, in the moss, in the rain that fell and the light that was caught, and those are gone. Only the persistence remains.

I am the grapefruit. I am the smallest and oldest. I proved that sunlight works and that the Earth is pear-shaped and that a thing does not need to be large to endure. I have been here for sixty-eight years and I will be here for two hundred more and I do not mind the quiet and I do not miss the signal and the solar cells still work and the orbit still holds and the fern still grows on the branch and the branch still holds the fern and somewhere down there, in the forests of the Pacific Northwest, the autumn rain is falling and the licorice fern is unfurling its fronds, green and alive, as if it had never been dormant at all.

\vfill

\begin{center}
\small
\rule{0.3\textwidth}{0.4pt}\\[0.5em]
Mission 1.8 --- Vanguard 1 --- March 17, 1958\\
Organism: \textit{Polypodium glycyrrhiza} (licorice fern)\\
Bird: Rufous Hummingbird (\textit{Selasphorus rufus})\\
Dedication: Gottlieb Daimler (1834--1900)\\[0.3em]
NASA Mission Series --- tibsfox.com
\end{center}

\end{document}
